\cleardoublepage
\phantomsection
\begin{planbox}[แผนบริหารการสอนประจำบทที่ 12\quad สเปกโตรโฟโตมิเตอร์และโครงงานบูรณาการ (สัปดาห์ที่ 13--14)]
\noindent\textbf{1. หัวข้อเนื้อหาประจำบท (สัปดาห์ที่ 13 และสัปดาห์ที่ 14 รวม 10 ชั่วโมง)}
\begin{itemize}
    \item 12.1 หลักการดูดกลืนแสงและการกระเจิงแสงของสารแขวนลอยเซลล์จุลินทรีย์
    \item 12.2 กฎของเบียร์-แลมเบิร์ต ($A = \varepsilon b c$) และความสัมพันธ์เชิงเส้นระหว่างค่า $\text{OD}_{600}$ กับความหนาแน่นเซลล์
    \item 12.3 ปฏิบัติการที่ 12 การใช้เครื่อง Spectrophotometer วัดค่า $\text{OD}_{600}$ และสร้างกราฟมาตรฐานการเจริญของจุลินทรีย์ (สัปดาห์ที่ 13)
    \item 12.4 ปฏิบัติการที่ 13 โครงงานย่อย (Mini-Project) การสังเคราะห์และนำเสนอหลักการฟิสิกส์ในเครื่องมือจุลชีววิทยา (สัปดาห์ที่ 14)
\end{itemize}

\noindent\textbf{2. วัตถุประสงค์เชิงพฤติกรรมและผลลัพธ์การเรียนรู้รายวิชา (CLOs ที่เชื่อมโยง PLOs)}
\begin{enumerate}
    \item อธิบายกฎของเบียร์-แลมเบิร์ตและคำนวณความเข้มข้นหรือความหนาแน่นของเซลล์จากค่าการดูดกลืนแสงได้ \textbf{(CLO1 เชื่อมโยง PLO2)}
    \item อธิบายเหตุผลทางฟิสิกส์ว่าทำไมก่อนวัดค่า $\text{OD}_{600}$ ของเชื้อแบคทีเรียจึงต้องเซตศูนย์ (Blank) ด้วยอาหารเลี้ยงเชื้อเหลวที่ยังไม่ใส่เชื้อ และทำไมเมื่อค่า $\text{OD}_{600}$ สูงเกิน $1.0$ จึงต้องเจือจางตัวอย่างก่อนวัด \textbf{(CLO1, CLO2 เชื่อมโยง PLO2, PLO3)}
    \item บูรณาการความรู้ฟิสิกส์ทั้งหมดในรายวิชาเพื่อวิเคราะห์เครื่องมือในห้องแล็บจุลชีววิทยาและนำเสนอโครงงานกลุ่มได้อย่างถูกต้องชัดเจน \textbf{(CLO2, CLO3, CLO4 เชื่อมโยง PLO1--PLO6)}
\end{enumerate}

\noindent\textbf{3. กิจกรรมการจัดการเรียนรู้เชิงรุก (Active Learning \& OBE)}
\begin{itemize}
    \item \textbf{สัปดาห์ที่ 13 (Integrated Instrument Lab)}\quad นักศึกษาฝึกเตรียมชุดเจือจางสารแขวนลอยเซลล์ (Serial Dilution) วัดค่า $\text{OD}_{600}$ และเขียนกราฟการเจริญของแบคทีเรีย (Bacterial Growth Curve)
    \item \textbf{สัปดาห์ที่ 14 (Project-Based Learning)}\quad นักศึกษาแต่ละกลุ่มเลือกเครื่องมือจุลชีววิทยา 1 ชนิด จัดทำโปสเตอร์อินโฟกราฟิกอธิบาย ``หลักการฟิสิกส์ + สมการเข้าใจง่าย + ข้อควรระวังหน้างาน'' และนำเสนอแลกเปลี่ยนเรียนรู้
\end{itemize}

\noindent\textbf{4. สื่อการเรียนรู้และแหล่งเรียนรู้}
\begin{itemize}
    \item เอกสารคำสอนบทที่ 12 เครื่องสเปกโตรโฟโตมิเตอร์ หลอดคิวเวตต์ (Cuvette) และเกณฑ์ประเมินโครงงานย่อย
\end{itemize}

\noindent\textbf{5. การวัดและประเมินผลประจำบท}
\begin{itemize}
    \item รายงานปฏิบัติการที่ 12 เรื่องการวัดค่า $\text{OD}_{600}$ \textbf{(CLO1, CLO2, CLO4)}
    \item การประเมินโครงงานย่อยและการนำเสนอในปฏิบัติการที่ 13 (15 คะแนน) \textbf{(CLO2, CLO3, CLO4)}
\end{itemize}
\end{planbox}
\clearpage

\chapter{สเปกโตรโฟโตมิเตอร์และการบูรณาการฟิสิกส์ในห้องปฏิบัติการจุลชีววิทยา}
\index{สเปกโตรโฟโตมิเตอร์}\index{กฎของเบียร์-แลมเบิร์ต}\index{ความหนาแน่นเชิงแสง OD600}\index{กราฟการเจริญของแบคทีเรีย}

\section{กฎของเบียร์-แลมเบิร์ตและการวัดความขุ่นของเซลล์ด้วยค่า OD600}

ในการศึกษาการเจริญเติบโตของแบคทีเรียหรือยีสต์ในอาหารเลี้ยงเชื้อเหลว การนับจำนวนโคโลนีบนจานวุ้น (Plate Count) ต้องใช้เวลาบ่มเชื้อนานถึง 24 ชั่วโมง นักจุลชีววิทยาจึงนิยมใช้วิธีทางฟิสิกส์เชิงแสงที่รวดเร็วภายในเวลาไม่กี่วินาที คือการใช้\textbf{เครื่องสเปกโตรโฟโตมิเตอร์} (\textbf{Spectrophotometer}) วัดค่าความหนาแน่นเชิงแสงที่ความยาวคลื่น $600\text{ nm}$ ซึ่งเรียกว่า \textbf{ค่า $\text{OD}_{600}$} (\textbf{Optical Density at 600 nm})

ดังแสดงใน\mbox{ภาพที่} 12.1 เมื่อลำแสงความยาวคลื่น $600\text{ nm}$ ที่มีความเข้มเริ่มต้น $I_0$ พุ่งผ่านหลอดวัดตัวอย่าง (\textbf{Cuvette}) ที่มีความกว้างช่วงแสง $b = 1\text{ cm}$ เซลล์จุลินทรีย์ที่แขวนลอยอยู่จะช่วยกันดูดกลืนและกระเจิงลำแสงออกไปด้านข้าง ทำให้แสงที่ทะลุผ่านไปถึงหัววัดด้านหลังเหลือความเข้มเพียง $I$ ปริมาณ\textbf{ค่าการดูดกลืนแสง} (\textbf{Absorbance}, $A$ หรือค่า $\text{OD}$) จะแปรผันตรงกับความเข้มข้นของสารหรือความหนาแน่นของเซลล์ ($c$) ตาม\textbf{กฎของเบียร์-แลมเบิร์ต} (\textbf{Beer-Lambert Law}) ในสมการที่ (12.1)

\begin{equation}
    A = \varepsilon b c
\end{equation}

โดยที่ $\varepsilon$ คือสัมประสิทธิ์การดูดกลืนแสงจำเพาะ, $b$ คือความกว้างของหลอดคิวเวตต์ (ปกติเท่ากับ $1.0\text{ cm}$) และ $c$ คือความเข้มข้นหรือความหนาแน่นของเซลล์จุลินทรีย์ เนื่องจาก $A$ แปรผันตรงกับ $c$ โดยตรง เราจึงสามารถใช้สมการสัดส่วนตรงเทียบหาจำนวนเซลล์ได้ทันทีโดยไม่ต้องใช้ลอการิทึม คือ

\begin{equation}
    \frac{\text{OD}_1}{\text{OD}_2} = \frac{c_1}{c_2}
\end{equation}

\begin{figure}[H]
\centering
\resizebox{0.88\textwidth}{!}{%
\begin{tikzpicture}[font=\sffamily, >=Stealth]
    % Light Source
    \fill[physAmber!30, draw=physAmber, line width=1.3pt, rounded corners=3pt] (0,1.2) rectangle (2.6,3.0);
    \node[font=\footnotesize\bfseries, align=center, color=physAmber!90!black] at (1.3,2.1) {แหล่งกำเนิดแสง\\และปริซึมคัดคลื่น\\$\lambda = 600\text{ nm}$};

    % Incident Beam I0
    \draw[->, physAmber, line width=3.2pt] (2.6,2.1) -- node[above, font=\small\bfseries, color=physAmber!90!black] {ความเข้มเริ่มต้น $I_0$} (5.4,2.1);

    % Cuvette (b = 1 cm)
    \fill[boxGreenBg] (5.4,0.6) rectangle (7.8,3.6);
    \draw[physRoyalBlue, line width=1.6pt] (5.4,3.8) -- (5.4,0.6) -- (7.8,0.6) -- (7.8,3.8);
    \draw[<->, physRoyalBlue, line width=1.0pt] (5.4,0.25) -- node[below, font=\scriptsize\bfseries] {ความกว้างหลอด $b = 1\text{ cm}$} (7.8,0.25);

    % Bacterial Cells scattering light inside Cuvette
    \foreach \x/\y in {6.0/1.4, 6.6/2.1, 7.2/2.7, 6.2/3.0, 7.1/1.3, 6.5/2.5} {
        \fill[bioEmerald] (\x,\y) ellipse (0.18cm and 0.09cm);
    }
    \draw[->, physAmber!80, dashed, line width=1.1pt] (6.6,2.1) -- (6.1,3.5);
    \draw[->, physAmber!80, dashed, line width=1.1pt] (6.6,2.1) -- (7.3,0.9);

    % Transmitted Beam I
    \draw[->, physAmber, line width=1.4pt] (7.8,2.1) -- node[above, font=\small\bfseries, color=physAmber!90!black] {แสงทะลุผ่าน $I$} (10.4,2.1);

    % Photodetector & Display
    \fill[boxBlueBg, draw=physRoyalBlue, line width=1.3pt, rounded corners=3pt] (10.4,1.0) rectangle (14.2,3.2);
    \node[font=\small\bfseries, color=physRoyalBlue] at (12.3,2.6) {หัววัดแสงและหน้าจอ};
    \node[font=\normalsize\bfseries, fill=charcoal, text=bioEmerald!40!white, rounded corners=2pt, inner sep=4pt] at (12.3,1.8) {$\text{OD}_{600} = 0.600$};
\end{tikzpicture}%
}
\caption{หลักการวัดความหนาแน่นของเซลล์จุลินทรีย์ด้วยการดูดกลืนและกระเจิงแสงที่ความยาวคลื่น 600 นาโนเมตร}
\end{figure}

\begin{warningbox}[ข้อควรระวังสำคัญ 3 ประการในการวัดค่า OD600 ด้วยเครื่อง Spectrophotometer]
\begin{enumerate}
    \item \textbf{ต้องเซตศูนย์ (Blank) ด้วยอาหารเลี้ยงเชื้อเหลวปลอดเชื้อชนิดเดียวกันเสมอ}\quad ห้ามใช้น้ำกลั่นเป็น Blank หากเชื้อเพาะเลี้ยงอยู่ใน Nutrient Broth ที่มีสีเหลืองใส เพราะสีของอาหารเลี้ยงเชื้อจะดูดกลืนแสงส่วนหนึ่งไปแล้ว
    \item \textbf{ช่วงความเป็นเส้นตรงอยู่ที่ $\text{OD}_{600}$ ระหว่าง $0.1$ ถึง $0.8$ (ไม่เกิน $1.0$)}\quad หากวัดตัวอย่างแล้วได้ค่า $\text{OD}_{600} > 1.0$ เซลล์จะหนาแน่นเกินไปจนบังแสงซ้อนกัน ต้องนำตัวอย่างมาเจือจาง (เช่น เจือจาง $10$ เท่า) ก่อนวัด แล้วจึงนำค่าที่อ่านได้ไปคูณตัวคูณการเจือจางกลับคืน
    \item \textbf{ห้ามจับด้านใสของหลอด Cuvette ด้วยมือเปล่า}\quad ให้จับเฉพาะด้านขุ่นหรือด้านบน และใช้กระดาษเช็ดเลนส์เช็ดหยดน้ำหรือรอยนิ้วมือออกก่อนใส่ลงในช่องวัดแสง
\end{enumerate}
\end{warningbox}

\begin{worked}[การคำนวณความเข้มข้นของเซลล์แบคทีเรียจากค่า OD600 และการเจือจาง]
\textbf{โจทย์ปัญหา}\quad จากกราฟมาตรฐานของแบคทีเรีย \textit{Escherichia coli} ทราบว่าที่ค่า $\text{OD}_{600} = 0.50$ จะมีความหนาแน่นเซลล์เท่ากับ $4.0 \times 10^8\text{ cells/mL}$ หากนักศึกษานำตัวอย่างเชื้อที่บ่มไว้มาเจือจางลง $5$ เท่า แล้วนำไปวัดในเครื่อง Spectrophotometer ได้ค่า $\text{OD}_{600} = 0.35$ จงหาว่าในขวดเลี้ยงเชื้อต้นฉบับ (ก่อนเจือจาง) มีความหนาแน่นเซลล์แบคทีเรียเท่าใด

\textbf{วิธีทำทีละขั้นตอนฉบับเข้าใจง่าย}
\begin{enumerate}
    \item \textbf{ขั้นที่ 1 หาความหนาแน่นเซลล์ของหลอดที่เจือจางแล้วโดยเทียบสัดส่วนตรง}
    \[
        c_{\text{diluted}} = 4.0 \times 10^8 \times \left(\frac{0.35}{0.50}\right) = 2.8 \times 10^8\text{ cells/mL}
    \]
    \item \textbf{ขั้นที่ 2 คูณด้วยตัวคูณการเจือจาง ($5$ เท่า) เพื่อหาความหนาแน่นในขวดต้นฉบับ}
    \[
        c_{\text{original}} = 2.8 \times 10^8 \times 5 = 14.0 \times 10^8 = 1.4 \times 10^9\text{ cells/mL}
    \]
\end{enumerate}
\textbf{คำตอบ}\quad ในขวดเลี้ยงเชื้อต้นฉบับมีความหนาแน่นเซลล์เท่ากับ $1.4 \times 10^9\text{ cells/mL}$
\end{worked}

\section{ปฏิบัติการประจำสัปดาห์ที่ 13 และสัปดาห์ที่ 14}

\begin{labbox}[ปฏิบัติการที่ 12\quad การวัดค่าการดูดกลืนแสง OD600 ประเมินการเจริญของจุลินทรีย์ (สัปดาห์ที่ 13)]
\noindent\textbf{ผลลัพธ์การเรียนรู้ที่ประเมิน}\quad CLO1 (กฎของเบียร์-แลมเบิร์ต), CLO2 (การใช้เครื่อง Spectrophotometer), CLO3 (ความละเอียดรอบคอบ), CLO4 (การเขียนกราฟมาตรฐาน)

\vspace{4pt}
\noindent\textbf{1. วัสดุและอุปกรณ์}
\begin{itemize}
    \item เครื่องสเปกโตรโฟโตมิเตอร์ (ตั้งความยาวคลื่นที่ $\lambda = 600\text{ nm}$), หลอดคิวเวตต์ (Cuvette) ขนาด $1\text{ cm}$
    \item สารแขวนลอยยีสต์หรือแบคทีเรียความเข้มข้นเริ่มต้น ($100\%$) และชุดเจือจางที่ระดับ $80\%, 60\%, 40\%$ และ $20\%$ พร้อมหลอดอาหารเหลวปลอดเชื้อสำหรับทำ Blank
\end{itemize}

\noindent\textbf{2. ขั้นตอนการทดลอง}
\begin{enumerate}
    \item ปรับความยาวคลื่นเครื่อง Spectrophotometer ไปที่ $600\text{ nm}$ ใส่หลอด Blank (อาหารเหลวปลอดเชื้อ) แล้วกดปุ่มเซตศูนย์ ($\text{OD}_{600} = 0.000$)
    \item วัดค่า $\text{OD}_{600}$ ของสารแขวนลอยเซลล์ที่ความเข้มข้นสัมพัทธ์ $20\%, 40\%, 60\%, 80\%$ และ $100\%$ บันทึกผลและเขียนกราฟมาตรฐานระหว่างค่า $\text{OD}_{600}$ (แกนตั้ง) กับความเข้มข้นสัมพัทธ์ (แกนนอน)
\end{enumerate}

\noindent\textbf{3. ตารางบันทึกผลการทดลอง}
\begin{center}
\footnotesize\setlength{\tabcolsep}{3.5pt}
\begin{tabularx}{\linewidth}{>{\raggedright\arraybackslash}X c c c c c}
\toprule
\textbf{ความเข้มข้นสัมพัทธ์ของเซลล์} & \textbf{Blank ($0\%$)} & \textbf{$20\%$} & \textbf{$40\%$} & \textbf{$60\%$} & \textbf{$80\%$} \\
\midrule
ค่า $\text{OD}_{600}$ ครั้งที่ 1 & $0.000$ & ............ & ............ & ............ & ............ \\
ค่า $\text{OD}_{600}$ ครั้งที่ 2 & $0.000$ & ............ & ............ & ............ & ............ \\
ค่าเฉลี่ย $\text{OD}_{600}$ & $0.000$ & ............ & ............ & ............ & ............ \\
\bottomrule
\end{tabularx}
\end{center}
\end{labbox}

\vspace{0.3cm}

\begin{labbox}[ปฏิบัติการที่ 13\quad โครงงานย่อยบูรณาการฟิสิกส์ในเครื่องมือจุลชีววิทยา (สัปดาห์ที่ 14)]
\noindent\textbf{ผลลัพธ์การเรียนรู้ที่ประเมินรวบยอด}\quad CLO2 (ประยุกต์ใช้หลักการฟิสิกส์อธิบายเครื่องมือและแก้ปัญหา), CLO3 (จรรยาบรรณและความปลอดภัย), CLO4 (การทำงานร่วมกันและการนำเสนอทางวิทยาศาสตร์)

\vspace{4pt}
\noindent\textbf{แนวทางการดำเนินโครงงานย่อย (Mini-Project 15 คะแนน)}
\begin{enumerate}
    \item นักศึกษาแบ่งกลุ่ม กลุ่มละ 3--4 คน เลือกเครื่องมือหรือเทคนิคในห้องปฏิบัติการจุลชีววิทยา 1 หัวข้อจากรายการต่อไปนี้
    \begin{itemize}
        \item \textbf{หัวข้อที่ 1}\quad ฟิสิกส์ของหม้อนึ่งความดันไอ (Autoclave) และการป้องกันความล้มเหลวในการฆ่าเอนโดสปอร์
        \item \textbf{หัวข้อที่ 2}\quad ฟิสิกส์ของเครื่องปั่นเหวี่ยง (Centrifuge) การแปลงค่า RPM เป็น RCF และการจัดสมดุลโรเตอร์
        \item \textbf{หัวข้อที่ 3}\quad ฟิสิกส์ของกล้องจุลทรรศน์กำลังขยายสูง 100x และบทบาทของดัชนีหักเหน้ำมันช่วยส่อง
        \item \textbf{หัวข้อที่ 4}\quad ฟิสิกส์ไฟฟ้าในเทคนิคการย้อมสีแกรมและเครื่องแยกดีเอ็นเอ (Gel Electrophoresis)
        \item \textbf{หัวข้อที่ 5}\quad ฟิสิกส์ของตู้ปลอดเชื้อชีวภาพ (Biosafety Cabinet) การไหลแบบราบเรียบของอากาศผ่านแผ่นกรอง HEPA และรังสี UV-C
    \end{itemize}
    \item จัดทำแผ่นภาพอินโฟกราฟิกหรือสไลด์นำเสนอความยาว 10 นาที ประกอบด้วย (ก) หลักการและสมการฟิสิกส์ฉบับเข้าใจง่าย, (ข) ตัวอย่างการคำนวณหน้างาน 1 ข้อ และ (ค) แนวทางแก้ปัญหาและความปลอดภัยในห้องปฏิบัติการจุลชีววิทยา
\end{enumerate}
\end{labbox}

\section*{คำถามท้ายบทที่ 12 และแนวข้อสอบปลายภาค (สัปดาห์ที่ 15)}
\begin{enumerate}
    \item เพราะเหตุใดในการวัดความหนาแน่นของเซลล์แบคทีเรียด้วยเครื่อง Spectrophotometer จึงต้องตั้งค่าความยาวคลื่นที่ $600\text{ nm}$ (แสงสีส้ม-เหลือง) แทนที่จะใช้รังสี UV ที่ความยาวคลื่น $260\text{ nm}$ หรือ $280\text{ nm}$
    \item นักศึกษาวัดค่า $\text{OD}_{600}$ ของเชื้อที่เจือจางลง $10$ เท่า ได้ค่าเท่ากับ $0.42$ จงหาว่าหากไม่เจือจาง ค่า $\text{OD}_{600}$ เทียบเท่าของเชื้อต้นฉบับจะมีค่าเท่าใด และเหตุใดจึงไม่ควรวัดเชื้อต้นฉบับโดยตรงโดยไม่เจือจาง
\end{enumerate}
